% =================================================================
%  CSE 815 -- Machine Learning
%  Module 3, Week 8, Part 2
%  Recommender Systems at Scale, and Principal Component Analysis
%  (the funnel and the PCA axis are step-by-step beamer overlays)
% =================================================================
\documentclass[aspectratio=169,10pt]{beamer}
\usepackage{cse815}

\title[CSE 815 -- Week 8, Part 2]{Recommender Systems at Scale, and PCA}
\subtitle{Related Items, Deep Content-Based Filtering, Large Catalogues, Ethics, and Principal Component Analysis}
\author{Rokan Uddin Faruqui\\ \small Professor, Department of CSE}
\institute{\small University of Chittagong \\ \texttt{rufaruqui@cu.ac.bd}}
\date[Week 8 / Part 2]{Week 8 $\cdot$ Part 2}

\begin{document}

\begin{frame}[plain,noframenumbering]
  \titlepage
\end{frame}

% ---------------------------------------------------------------
\begin{frame}{From a toy matrix to millions of items}
  \begin{columns}[T,onlytextwidth]
    \begin{column}{0.5\textwidth}
      \begin{block}{Part 1 gave us}
        \footnotesize
        \begin{itemize}
          \item Content-based prediction: $\vw^{(j)}\cdot\vx^{(i)} + b^{(j)}$ with known
                item features.
          \item Collaborative filtering: learn $\vx$, $\vw$, $b$ together from the
                ratings alone.
          \item Binary labels, mean normalization, a \texttt{GradientTape} loop.
        \end{itemize}
      \end{block}
      \footnotesize
      Real systems have millions of users and items, rich side information about both,
      and a few hundred milliseconds to answer. Today: how they cope --- and a
      dimensionality-reduction method that grew out of the same idea of compact
      representations.
    \end{column}
    \begin{column}{0.46\textwidth}
      \begin{keyidea}[Today]
        \footnotesize
        \begin{enumerate}
          \item Finding related items
          \item Collaborative vs.\ content-based filtering
          \item Deep learning for content-based filtering
          \item Recommending from a large catalogue
          \item Ethical use of recommender systems
          \item Principal component analysis (optional)
        \end{enumerate}
      \end{keyidea}
    \end{column}
  \end{columns}
\end{frame}

% ===============================================================
\section{Finding Related Items}

\begin{frame}{``Because you ordered \emph{Pepper Soup}'' \dots}
  \begin{columns}[T,onlytextwidth]
    \begin{column}{0.5\textwidth}
      \begin{block}{Use the learned features}
        \footnotesize
        Collaborative filtering gave every dish a feature vector $\vx^{(i)}$. Dishes
        liked by the same people end up close together, so the items most related to
        dish $i$ are those minimising
        \[ \bigl\|\vx^{(k)} - \vx^{(i)}\bigr\|^2 = \sum_{l=1}^{n}\bigl(x^{(k)}_l - x^{(i)}_l\bigr)^2 . \]
      \end{block}
      \begin{keyidea}
        No cuisine labels, no recipes: relatedness comes entirely from co-ratings. This
        is ``customers who liked this also liked'' made precise.
      \end{keyidea}
    \end{column}
    \begin{column}{0.46\textwidth}
      \footnotesize
      \begin{block}{Distances from \emph{Pepper Soup}}
        \renewcommand{\arraystretch}{1.18}
        \begin{tabular}{@{}lr@{}}
          \toprule
          dish & $\|\vx^{(k)} - \vx^{(3)}\|^2$\\
          \midrule
        Chilli Noodles & $0.091$\\
        Chicken Curry & $0.856$\\
        Mango Ice Cream & $7.573$\\
        Ginger Cake & $7.617$\\
          \bottomrule
        \end{tabular}
      \end{block}
      Computed from the features learned in Part 1's animation. \emph{Chilli Noodles}
      is nearest; the sweet dishes are more than eighty times as far away.
    \end{column}
  \end{columns}
\end{frame}

% ===============================================================
\section{Collaborative vs.\ Content-Based Filtering}

\begin{frame}{Two sources of information}
  \footnotesize
  \renewcommand{\arraystretch}{1.3}
  \begin{tabular}{@{}L{2.8cm}L{5.2cm}L{5.2cm}@{}}
    \toprule
    & \textbf{collaborative filtering} & \textbf{content-based filtering}\\
    \midrule
    uses & ratings of users who rated similarly & features of the user and of the item\\
    item features & learned from ratings & given: cuisine, price, spice level, ingredients, photos\\
    new item & cannot be recommended until rated (\emph{cold start}) & recommended at once from its features\\
    new user & needs a few ratings (or mean normalization) & needs profile features (age, country, stated interests)\\
    strength & discovers tastes no feature describes & works from day one; uses side information\\
    \bottomrule
  \end{tabular}
  \vspace{3mm}
  \begin{columns}[T,onlytextwidth]
    \begin{column}{0.5\textwidth}
      \begin{block}{Features for content-based filtering}
        \textbf{User} $\vx_u$: age, location, past ratings per cuisine, dishes ordered.\\
        \textbf{Dish} $\vx_m$: cuisine, price, spice level, reviews, average rating.\\
        The two vectors can have different lengths.
      \end{block}
    \end{column}
    \begin{column}{0.46\textwidth}
      \begin{keyidea}
        Large systems use both: content features for new items and new users,
        collaborative signals once interaction data accumulates.
      \end{keyidea}
    \end{column}
  \end{columns}
\end{frame}

% ===============================================================
\section{Deep Learning for Content-Based Filtering}

\begin{frame}{Two towers}
  \begin{columns}[T,onlytextwidth]
    \begin{column}{0.55\textwidth}
      \centering
      \begin{tikzpicture}[
          box/.style={draw=cublue, fill=culight, rounded corners, align=center,
                      font=\scriptsize, minimum width=1.9cm, minimum height=7mm},
          ar/.style={-{Stealth[length=1.8mm]}, cugray, thick}]
        \node[box, draw=cuteal, fill=cuteal!12] (xu) at (0,0) {user features\\$\vx_u$};
        \node[box] (nu) at (0,-1.3) {user network\\$128\to64\to32$};
        \node[box, draw=cuamber, fill=cuamber!15] (vu) at (0,-2.6) {$\mathbf{v}_u\in\R^{32}$};
        \node[box, draw=cuteal, fill=cuteal!12] (xm) at (3.4,0) {dish features\\$\vx_m$};
        \node[box] (nm) at (3.4,-1.3) {dish network\\$128\to64\to32$};
        \node[box, draw=cuamber, fill=cuamber!15] (vm) at (3.4,-2.6) {$\mathbf{v}_m\in\R^{32}$};
        \node[box, draw=curose, fill=curose!8, minimum width=3.2cm] (dot) at (1.7,-3.9)
          {prediction $\mathbf{v}_u\cdot\mathbf{v}_m$\\(or $g(\mathbf{v}_u\cdot\mathbf{v}_m)$ for likes)};
        \draw[ar] (xu) -- (nu); \draw[ar] (nu) -- (vu);
        \draw[ar] (xm) -- (nm); \draw[ar] (nm) -- (vm);
        \draw[ar] (vu) -- (dot); \draw[ar] (vm) -- (dot);
      \end{tikzpicture}
    \end{column}
    \begin{column}{0.41\textwidth}
      \footnotesize
      \begin{block}{The idea}
        One network maps user features to $\mathbf{v}_u$, another maps dish features
        to $\mathbf{v}_m$, of the \emph{same} length. Predict with their dot product ---
        Part 1's $\vw^{(j)}\cdot\vx^{(i)}$, both sides now computed.
      \end{block}
      \begin{block}{Training}
        Both towers are trained together on one cost:
        \[ J = \!\!\sum_{r(i,j)=1}\!\!\bigl(\mathbf{v}_u^{(j)}\cdot\mathbf{v}_m^{(i)} - y^{(i,j)}\bigr)^2 + \text{penalty} \]
      \end{block}
      Related dishes: small $\|\mathbf{v}_m^{(k)} - \mathbf{v}_m^{(i)}\|^2$,
      precomputed overnight for the whole catalogue.
    \end{column}
  \end{columns}
\end{frame}

\begin{frame}[fragile]{Two towers in TensorFlow}
  \begin{columns}[T,onlytextwidth]
    \begin{column}{0.56\textwidth}
\begin{lstlisting}
user_NN = Sequential([Dense(256, activation='relu'),
                      Dense(128, activation='relu'),
                      Dense(32)])
item_NN = Sequential([Dense(256, activation='relu'),
                      Dense(128, activation='relu'),
                      Dense(32)])
u_in = Input(shape=(n_user_feats,))
m_in = Input(shape=(n_item_feats,))
vu = tf.linalg.l2_normalize(user_NN(u_in), axis=1)
vm = tf.linalg.l2_normalize(item_NN(m_in), axis=1)
out = Dot(axes=1)([vu, vm])
model = Model([u_in, m_in], out)
model.compile(optimizer=Adam(), loss=MeanSquaredError())
\end{lstlisting}
    \end{column}
    \begin{column}{0.4\textwidth}
      \footnotesize
      \begin{block}{What is new}
        Two \texttt{Sequential} networks, combined with the functional API:
        \texttt{Model} takes two inputs and one output. \texttt{Dot} computes
        $\mathbf{v}_u\cdot\mathbf{v}_m$ for each pair.
      \end{block}
      \begin{block}{Normalizing the vectors}
        \texttt{l2\_normalize} gives every vector length $1$, so the dot product is the
        cosine of the angle between them --- a scale-free similarity that tends to train
        more stably.
      \end{block}
    \end{column}
  \end{columns}
\end{frame}

% ===============================================================
\section{Recommending from a Large Catalogue}

\begin{frame}{Retrieval and ranking}
  \begin{columns}[T,onlytextwidth]
    \begin{column}{0.55\textwidth}
      \centering
      \begin{tikzpicture}[stage/.style={draw, rounded corners, thick, align=center,
          font=\scriptsize, minimum height=9mm}]
        \node[stage, draw=cublue, fill=culight, minimum width=6.2cm] (c) at (0,0)
          {catalogue: millions of items};
        \onslide<2->{
        \node[stage, draw=cuteal, fill=cuteal!12, minimum width=4.6cm] (r) at (0,-1.35)
          {\textbf{retrieval}: a few hundred plausible candidates};
        \draw[-{Stealth}, thick, cugray] (c) -- (r);}
        \onslide<3->{
        \node[stage, draw=cuamber, fill=cuamber!15, minimum width=3.2cm] (k) at (0,-2.7)
          {\textbf{ranking}: score each\\with the full model};
        \draw[-{Stealth}, thick, cugray] (r) -- (k);}
        \onslide<4->{
        \node[stage, draw=curose, fill=curose!8, minimum width=2.0cm] (s) at (0,-4.0)
          {show top 10};
        \draw[-{Stealth}, thick, cugray] (k) -- (s);}
      \end{tikzpicture}
    \end{column}
    \begin{column}{0.41\textwidth}
      \footnotesize
      \only<1>{\begin{block}{The problem}
        Scoring every one of millions of items with a deep network, for every user
        request, in a fraction of a second, is too slow and too expensive.
      \end{block}}
      \only<2>{\begin{block}{Step 1: retrieval}
        Cheaply generate a broad candidate list: the $10$ items most similar to each of
        the user's last $10$ ordered; the most popular dishes in their three favourite
        cuisines; top dishes in their area. Remove duplicates and restaurants that
        are closed or too far away.
      \end{block}}
      \only<3>{\begin{block}{Step 2: ranking}
        Run the two-tower model on just these candidates. $\mathbf{v}_m$ for every item
        is precomputed; only $\mathbf{v}_u$ is computed per request --- one network
        evaluation, then a few hundred dot products.
      \end{block}}
      \only<4>{\begin{block}{Step 3: display}
        Show the highest-scoring items. Retrieving more candidates improves the
        recommendations but slows them down: tune the number retrieved with offline
        experiments that measure the trade-off.
      \end{block}}
      \begin{keyidea}
        Retrieval must be fast and high-recall; ranking must be accurate. Splitting the
        two lets each be optimised for its own job.
      \end{keyidea}
    \end{column}
  \end{columns}
\end{frame}

% ===============================================================
\section{Ethical Use of Recommender Systems}

\begin{frame}{What should a recommender optimise?}
  \begin{columns}[T,onlytextwidth]
    \begin{column}{0.5\textwidth}
      \footnotesize
      \begin{block}{Possible goals}
        \begin{itemize}
          \item items the user is most likely to rate highly;
          \item items most likely to be bought;
          \item advertisements most likely to be clicked --- and paid for;
          \item items that generate the most profit;
          \item videos that maximise watch time.
        \end{itemize}
      \end{block}
      \begin{block}{When the goal diverges from the user's good}
        \begin{itemize}
          \item Watch-time maximisation can amplify outrage, conspiracy and extreme
                content, because it engages.
          \item Advertising auctions can reward the businesses with the highest margins,
                not the best service.
          \item Feedback loops: what is recommended gets clicked, so gets recommended
                more, narrowing what users ever see.
        \end{itemize}
      \end{block}
    \end{column}
    \begin{column}{0.46\textwidth}
      \begin{alertblock}{Choosing a metric is choosing a value}
        \footnotesize ``Maximise engagement'' sounds neutral. It is a decision that the
        system should show people whatever holds their attention, whether or not they
        are better off for it.
      \end{alertblock}
      \footnotesize
      \begin{block}{Practices that help}
        \begin{itemize}
          \item Be transparent with users about why an item is shown.
          \item Filter out problematic content and exploitative sellers.
          \item Measure long-term satisfaction, not only clicks.
          \item Inject diversity; audit who and what is never recommended.
        \end{itemize}
      \end{block}
    \end{column}
  \end{columns}
\end{frame}

% ===============================================================
\section{Principal Component Analysis}

\begin{frame}{Reducing the number of features}
  \begin{columns}[T,onlytextwidth]
    \begin{column}{0.5\textwidth}
      \footnotesize
      \begin{block}{Why}
        Data with $50$ or $1000$ features cannot be plotted. Often the features are
        strongly correlated --- a car's length and width, a country's GDP and its
        energy use --- so the data lies close to a lower-dimensional surface.
      \end{block}
      \begin{defnbox}[Principal component analysis (PCA)]
        \small Find new axes --- the \textbf{principal components} --- such that projecting
        the data onto the first $k$ of them keeps as much of its \emph{variance} (spread)
        as possible. Replace each $\vx\in\R^n$ by its $k$ coordinates
        $z_1,\dots,z_k$, with $k \ll n$.
      \end{defnbox}
    \end{column}
    \begin{column}{0.46\textwidth}
      \footnotesize
      \begin{block}{Uses}
        \begin{itemize}
          \item \textbf{Visualisation}: reduce to $2$ or $3$ dimensions and plot.
          \item \textbf{Compression}: store $k$ numbers instead of $n$.
          \item Faster training (less common now).
        \end{itemize}
      \end{block}
      \begin{keyidea}
        Like collaborative filtering, PCA finds a compact representation in which a few
        learned numbers summarise many observed ones. Unlike it, PCA is unsupervised
        and has an exact solution.
      \end{keyidea}
    \end{column}
  \end{columns}
\end{frame}

\begin{frame}{Searching for the best axis}
  \begin{columns}[T,onlytextwidth]
    \begin{column}{0.52\textwidth}
      \centering
      \begin{tikzpicture}
        \begin{axis}[width=6.4cm,height=6.0cm,xmin=-3,xmax=3,ymin=-3,ymax=3,
          axis equal image,xlabel={$x_1$ (standardized)},ylabel={$x_2$},
          tick label style={font=\tiny},label style={font=\scriptsize}]
          \addplot[only marks, mark=*, mark size=1.4pt, cublue] coordinates {(-1.726,-1.542)(-1.799,-1.496)(-1.138,-0.810)(0.261,-0.470)(-2.136,-1.390)(0.438,-0.332)(-0.255,-1.012)(0.653,0.592)(0.203,1.050)(0.986,1.264)(1.187,1.170)(-0.188,-0.293)(0.417,0.835)(1.056,2.038)(-0.473,-0.701)(0.661,0.709)(0.355,0.534)(1.733,1.447)(-0.933,-1.014)(-0.204,-1.093)(0.087,1.207)(0.374,0.360)(-1.728,-0.991)(0.064,0.100)(0.683,0.563)(-0.088,-0.044)(0.181,-0.151)(1.119,-0.188)(-1.275,-1.490)(1.485,1.146)};
\only<1>{\addplot[cuamber, very thick, domain=-3:3, samples=2] ({x*1.0000},{x*0.0000});
\draw[cugray!60] (axis cs:-1.726,-1.542) -- (axis cs:-1.726,-0.000); \draw[cugray!60] (axis cs:-1.799,-1.496) -- (axis cs:-1.799,-0.000); \draw[cugray!60] (axis cs:-1.138,-0.810) -- (axis cs:-1.138,-0.000); \draw[cugray!60] (axis cs:0.261,-0.470) -- (axis cs:0.261,0.000); \draw[cugray!60] (axis cs:-2.136,-1.390) -- (axis cs:-2.136,-0.000); \draw[cugray!60] (axis cs:0.438,-0.332) -- (axis cs:0.438,0.000); \draw[cugray!60] (axis cs:-0.255,-1.012) -- (axis cs:-0.255,-0.000); \draw[cugray!60] (axis cs:0.653,0.592) -- (axis cs:0.653,0.000); \draw[cugray!60] (axis cs:0.203,1.050) -- (axis cs:0.203,0.000); \draw[cugray!60] (axis cs:0.986,1.264) -- (axis cs:0.986,0.000); \draw[cugray!60] (axis cs:1.187,1.170) -- (axis cs:1.187,0.000); \draw[cugray!60] (axis cs:-0.188,-0.293) -- (axis cs:-0.188,-0.000); \draw[cugray!60] (axis cs:0.417,0.835) -- (axis cs:0.417,0.000); \draw[cugray!60] (axis cs:1.056,2.038) -- (axis cs:1.056,0.000); \draw[cugray!60] (axis cs:-0.473,-0.701) -- (axis cs:-0.473,-0.000); \draw[cugray!60] (axis cs:0.661,0.709) -- (axis cs:0.661,0.000); \draw[cugray!60] (axis cs:0.355,0.534) -- (axis cs:0.355,0.000); \draw[cugray!60] (axis cs:1.733,1.447) -- (axis cs:1.733,0.000); \draw[cugray!60] (axis cs:-0.933,-1.014) -- (axis cs:-0.933,-0.000); \draw[cugray!60] (axis cs:-0.204,-1.093) -- (axis cs:-0.204,-0.000); \draw[cugray!60] (axis cs:0.087,1.207) -- (axis cs:0.087,0.000); \draw[cugray!60] (axis cs:0.374,0.360) -- (axis cs:0.374,0.000); \draw[cugray!60] (axis cs:-1.728,-0.991) -- (axis cs:-1.728,-0.000); \draw[cugray!60] (axis cs:0.064,0.100) -- (axis cs:0.064,0.000); \draw[cugray!60] (axis cs:0.683,0.563) -- (axis cs:0.683,0.000); \draw[cugray!60] (axis cs:-0.088,-0.044) -- (axis cs:-0.088,-0.000); \draw[cugray!60] (axis cs:0.181,-0.151) -- (axis cs:0.181,0.000); \draw[cugray!60] (axis cs:1.119,-0.188) -- (axis cs:1.119,0.000); \draw[cugray!60] (axis cs:-1.275,-1.490) -- (axis cs:-1.275,-0.000); \draw[cugray!60] (axis cs:1.485,1.146) -- (axis cs:1.485,0.000);
\addplot[only marks, mark=*, mark size=1.1pt, cuamber] coordinates {(-1.726,-0.000)(-1.799,-0.000)(-1.138,-0.000)(0.261,0.000)(-2.136,-0.000)(0.438,0.000)(-0.255,-0.000)(0.653,0.000)(0.203,0.000)(0.986,0.000)(1.187,0.000)(-0.188,-0.000)(0.417,0.000)(1.056,0.000)(-0.473,-0.000)(0.661,0.000)(0.355,0.000)(1.733,0.000)(-0.933,-0.000)(-0.204,-0.000)(0.087,0.000)(0.374,0.000)(-1.728,-0.000)(0.064,0.000)(0.683,0.000)(-0.088,-0.000)(0.181,0.000)(1.119,0.000)(-1.275,-0.000)(1.485,0.000)};}%
\only<2>{\addplot[cuamber, very thick, domain=-3:3, samples=2] ({x*0.9659},{x*0.2588});
\draw[cugray!60] (axis cs:-1.726,-1.542) -- (axis cs:-1.996,-0.535); \draw[cugray!60] (axis cs:-1.799,-1.496) -- (axis cs:-2.053,-0.550); \draw[cugray!60] (axis cs:-1.138,-0.810) -- (axis cs:-1.264,-0.339); \draw[cugray!60] (axis cs:0.261,-0.470) -- (axis cs:0.126,0.034); \draw[cugray!60] (axis cs:-2.136,-1.390) -- (axis cs:-2.341,-0.627); \draw[cugray!60] (axis cs:0.438,-0.332) -- (axis cs:0.326,0.087); \draw[cugray!60] (axis cs:-0.255,-1.012) -- (axis cs:-0.491,-0.131); \draw[cugray!60] (axis cs:0.653,0.592) -- (axis cs:0.757,0.203); \draw[cugray!60] (axis cs:0.203,1.050) -- (axis cs:0.452,0.121); \draw[cugray!60] (axis cs:0.986,1.264) -- (axis cs:1.236,0.331); \draw[cugray!60] (axis cs:1.187,1.170) -- (axis cs:1.400,0.375); \draw[cugray!60] (axis cs:-0.188,-0.293) -- (axis cs:-0.249,-0.067); \draw[cugray!60] (axis cs:0.417,0.835) -- (axis cs:0.598,0.160); \draw[cugray!60] (axis cs:1.056,2.038) -- (axis cs:1.494,0.400); \draw[cugray!60] (axis cs:-0.473,-0.701) -- (axis cs:-0.616,-0.165); \draw[cugray!60] (axis cs:0.661,0.709) -- (axis cs:0.794,0.213); \draw[cugray!60] (axis cs:0.355,0.534) -- (axis cs:0.465,0.125); \draw[cugray!60] (axis cs:1.733,1.447) -- (axis cs:1.978,0.530); \draw[cugray!60] (axis cs:-0.933,-1.014) -- (axis cs:-1.124,-0.301); \draw[cugray!60] (axis cs:-0.204,-1.093) -- (axis cs:-0.464,-0.124); \draw[cugray!60] (axis cs:0.087,1.207) -- (axis cs:0.383,0.103); \draw[cugray!60] (axis cs:0.374,0.360) -- (axis cs:0.439,0.118); \draw[cugray!60] (axis cs:-1.728,-0.991) -- (axis cs:-1.860,-0.498); \draw[cugray!60] (axis cs:0.064,0.100) -- (axis cs:0.085,0.023); \draw[cugray!60] (axis cs:0.683,0.563) -- (axis cs:0.778,0.208); \draw[cugray!60] (axis cs:-0.088,-0.044) -- (axis cs:-0.094,-0.025); \draw[cugray!60] (axis cs:0.181,-0.151) -- (axis cs:0.131,0.035); \draw[cugray!60] (axis cs:1.119,-0.188) -- (axis cs:0.997,0.267); \draw[cugray!60] (axis cs:-1.275,-1.490) -- (axis cs:-1.562,-0.419); \draw[cugray!60] (axis cs:1.485,1.146) -- (axis cs:1.672,0.448);
\addplot[only marks, mark=*, mark size=1.1pt, cuamber] coordinates {(-1.996,-0.535)(-2.053,-0.550)(-1.264,-0.339)(0.126,0.034)(-2.341,-0.627)(0.326,0.087)(-0.491,-0.131)(0.757,0.203)(0.452,0.121)(1.236,0.331)(1.400,0.375)(-0.249,-0.067)(0.598,0.160)(1.494,0.400)(-0.616,-0.165)(0.794,0.213)(0.465,0.125)(1.978,0.530)(-1.124,-0.301)(-0.464,-0.124)(0.383,0.103)(0.439,0.118)(-1.860,-0.498)(0.085,0.023)(0.778,0.208)(-0.094,-0.025)(0.131,0.035)(0.997,0.267)(-1.562,-0.419)(1.672,0.448)};}%
\only<3>{\addplot[cuamber, very thick, domain=-3:3, samples=2] ({x*0.8660},{x*0.5000});
\draw[cugray!60] (axis cs:-1.726,-1.542) -- (axis cs:-1.962,-1.133); \draw[cugray!60] (axis cs:-1.799,-1.496) -- (axis cs:-1.997,-1.153); \draw[cugray!60] (axis cs:-1.138,-0.810) -- (axis cs:-1.204,-0.695); \draw[cugray!60] (axis cs:0.261,-0.470) -- (axis cs:-0.007,-0.004); \draw[cugray!60] (axis cs:-2.136,-1.390) -- (axis cs:-2.204,-1.272); \draw[cugray!60] (axis cs:0.438,-0.332) -- (axis cs:0.185,0.107); \draw[cugray!60] (axis cs:-0.255,-1.012) -- (axis cs:-0.629,-0.363); \draw[cugray!60] (axis cs:0.653,0.592) -- (axis cs:0.746,0.431); \draw[cugray!60] (axis cs:0.203,1.050) -- (axis cs:0.607,0.350); \draw[cugray!60] (axis cs:0.986,1.264) -- (axis cs:1.287,0.743); \draw[cugray!60] (axis cs:1.187,1.170) -- (axis cs:1.397,0.807); \draw[cugray!60] (axis cs:-0.188,-0.293) -- (axis cs:-0.268,-0.155); \draw[cugray!60] (axis cs:0.417,0.835) -- (axis cs:0.675,0.390); \draw[cugray!60] (axis cs:1.056,2.038) -- (axis cs:1.674,0.967); \draw[cugray!60] (axis cs:-0.473,-0.701) -- (axis cs:-0.658,-0.380); \draw[cugray!60] (axis cs:0.661,0.709) -- (axis cs:0.803,0.464); \draw[cugray!60] (axis cs:0.355,0.534) -- (axis cs:0.498,0.287); \draw[cugray!60] (axis cs:1.733,1.447) -- (axis cs:1.926,1.112); \draw[cugray!60] (axis cs:-0.933,-1.014) -- (axis cs:-1.139,-0.658); \draw[cugray!60] (axis cs:-0.204,-1.093) -- (axis cs:-0.626,-0.362); \draw[cugray!60] (axis cs:0.087,1.207) -- (axis cs:0.588,0.339); \draw[cugray!60] (axis cs:0.374,0.360) -- (axis cs:0.436,0.252); \draw[cugray!60] (axis cs:-1.728,-0.991) -- (axis cs:-1.725,-0.996); \draw[cugray!60] (axis cs:0.064,0.100) -- (axis cs:0.091,0.053); \draw[cugray!60] (axis cs:0.683,0.563) -- (axis cs:0.756,0.436); \draw[cugray!60] (axis cs:-0.088,-0.044) -- (axis cs:-0.085,-0.049); \draw[cugray!60] (axis cs:0.181,-0.151) -- (axis cs:0.071,0.041); \draw[cugray!60] (axis cs:1.119,-0.188) -- (axis cs:0.758,0.438); \draw[cugray!60] (axis cs:-1.275,-1.490) -- (axis cs:-1.602,-0.925); \draw[cugray!60] (axis cs:1.485,1.146) -- (axis cs:1.610,0.930);
\addplot[only marks, mark=*, mark size=1.1pt, cuamber] coordinates {(-1.962,-1.133)(-1.997,-1.153)(-1.204,-0.695)(-0.007,-0.004)(-2.204,-1.272)(0.185,0.107)(-0.629,-0.363)(0.746,0.431)(0.607,0.350)(1.287,0.743)(1.397,0.807)(-0.268,-0.155)(0.675,0.390)(1.674,0.967)(-0.658,-0.380)(0.803,0.464)(0.498,0.287)(1.926,1.112)(-1.139,-0.658)(-0.626,-0.362)(0.588,0.339)(0.436,0.252)(-1.725,-0.996)(0.091,0.053)(0.756,0.436)(-0.085,-0.049)(0.071,0.041)(0.758,0.438)(-1.602,-0.925)(1.610,0.930)};}%
\only<4>{\addplot[cuteal, very thick, domain=-3:3, samples=2] ({x*0.7071},{x*0.7071});
\draw[cugray!60] (axis cs:-1.726,-1.542) -- (axis cs:-1.634,-1.634); \draw[cugray!60] (axis cs:-1.799,-1.496) -- (axis cs:-1.648,-1.648); \draw[cugray!60] (axis cs:-1.138,-0.810) -- (axis cs:-0.974,-0.974); \draw[cugray!60] (axis cs:0.261,-0.470) -- (axis cs:-0.104,-0.104); \draw[cugray!60] (axis cs:-2.136,-1.390) -- (axis cs:-1.763,-1.763); \draw[cugray!60] (axis cs:0.438,-0.332) -- (axis cs:0.053,0.053); \draw[cugray!60] (axis cs:-0.255,-1.012) -- (axis cs:-0.633,-0.633); \draw[cugray!60] (axis cs:0.653,0.592) -- (axis cs:0.622,0.622); \draw[cugray!60] (axis cs:0.203,1.050) -- (axis cs:0.626,0.626); \draw[cugray!60] (axis cs:0.986,1.264) -- (axis cs:1.125,1.125); \draw[cugray!60] (axis cs:1.187,1.170) -- (axis cs:1.179,1.179); \draw[cugray!60] (axis cs:-0.188,-0.293) -- (axis cs:-0.241,-0.241); \draw[cugray!60] (axis cs:0.417,0.835) -- (axis cs:0.626,0.626); \draw[cugray!60] (axis cs:1.056,2.038) -- (axis cs:1.547,1.547); \draw[cugray!60] (axis cs:-0.473,-0.701) -- (axis cs:-0.587,-0.587); \draw[cugray!60] (axis cs:0.661,0.709) -- (axis cs:0.685,0.685); \draw[cugray!60] (axis cs:0.355,0.534) -- (axis cs:0.445,0.445); \draw[cugray!60] (axis cs:1.733,1.447) -- (axis cs:1.590,1.590); \draw[cugray!60] (axis cs:-0.933,-1.014) -- (axis cs:-0.974,-0.974); \draw[cugray!60] (axis cs:-0.204,-1.093) -- (axis cs:-0.648,-0.648); \draw[cugray!60] (axis cs:0.087,1.207) -- (axis cs:0.647,0.647); \draw[cugray!60] (axis cs:0.374,0.360) -- (axis cs:0.367,0.367); \draw[cugray!60] (axis cs:-1.728,-0.991) -- (axis cs:-1.359,-1.359); \draw[cugray!60] (axis cs:0.064,0.100) -- (axis cs:0.082,0.082); \draw[cugray!60] (axis cs:0.683,0.563) -- (axis cs:0.623,0.623); \draw[cugray!60] (axis cs:-0.088,-0.044) -- (axis cs:-0.066,-0.066); \draw[cugray!60] (axis cs:0.181,-0.151) -- (axis cs:0.015,0.015); \draw[cugray!60] (axis cs:1.119,-0.188) -- (axis cs:0.466,0.466); \draw[cugray!60] (axis cs:-1.275,-1.490) -- (axis cs:-1.383,-1.383); \draw[cugray!60] (axis cs:1.485,1.146) -- (axis cs:1.315,1.315);
\addplot[only marks, mark=*, mark size=1.1pt, cuteal] coordinates {(-1.634,-1.634)(-1.648,-1.648)(-0.974,-0.974)(-0.104,-0.104)(-1.763,-1.763)(0.053,0.053)(-0.633,-0.633)(0.622,0.622)(0.626,0.626)(1.125,1.125)(1.179,1.179)(-0.241,-0.241)(0.626,0.626)(1.547,1.547)(-0.587,-0.587)(0.685,0.685)(0.445,0.445)(1.590,1.590)(-0.974,-0.974)(-0.648,-0.648)(0.647,0.647)(0.367,0.367)(-1.359,-1.359)(0.082,0.082)(0.623,0.623)(-0.066,-0.066)(0.015,0.015)(0.466,0.466)(-1.383,-1.383)(1.315,1.315)};}%
\only<5>{\addplot[cuamber, very thick, domain=-3:3, samples=2] ({x*0.5000},{x*0.8660});
\draw[cugray!60] (axis cs:-1.726,-1.542) -- (axis cs:-1.099,-1.904); \draw[cugray!60] (axis cs:-1.799,-1.496) -- (axis cs:-1.098,-1.901); \draw[cugray!60] (axis cs:-1.138,-0.810) -- (axis cs:-0.635,-1.100); \draw[cugray!60] (axis cs:0.261,-0.470) -- (axis cs:-0.138,-0.239); \draw[cugray!60] (axis cs:-2.136,-1.390) -- (axis cs:-1.136,-1.967); \draw[cugray!60] (axis cs:0.438,-0.332) -- (axis cs:-0.034,-0.059); \draw[cugray!60] (axis cs:-0.255,-1.012) -- (axis cs:-0.502,-0.869); \draw[cugray!60] (axis cs:0.653,0.592) -- (axis cs:0.419,0.727); \draw[cugray!60] (axis cs:0.203,1.050) -- (axis cs:0.505,0.875); \draw[cugray!60] (axis cs:0.986,1.264) -- (axis cs:0.794,1.375); \draw[cugray!60] (axis cs:1.187,1.170) -- (axis cs:0.804,1.392); \draw[cugray!60] (axis cs:-0.188,-0.293) -- (axis cs:-0.174,-0.301); \draw[cugray!60] (axis cs:0.417,0.835) -- (axis cs:0.466,0.807); \draw[cugray!60] (axis cs:1.056,2.038) -- (axis cs:1.146,1.986); \draw[cugray!60] (axis cs:-0.473,-0.701) -- (axis cs:-0.422,-0.730); \draw[cugray!60] (axis cs:0.661,0.709) -- (axis cs:0.472,0.818); \draw[cugray!60] (axis cs:0.355,0.534) -- (axis cs:0.320,0.555); \draw[cugray!60] (axis cs:1.733,1.447) -- (axis cs:1.060,1.836); \draw[cugray!60] (axis cs:-0.933,-1.014) -- (axis cs:-0.672,-1.165); \draw[cugray!60] (axis cs:-0.204,-1.093) -- (axis cs:-0.524,-0.908); \draw[cugray!60] (axis cs:0.087,1.207) -- (axis cs:0.544,0.943); \draw[cugray!60] (axis cs:0.374,0.360) -- (axis cs:0.249,0.432); \draw[cugray!60] (axis cs:-1.728,-0.991) -- (axis cs:-0.861,-1.491); \draw[cugray!60] (axis cs:0.064,0.100) -- (axis cs:0.059,0.103); \draw[cugray!60] (axis cs:0.683,0.563) -- (axis cs:0.414,0.718); \draw[cugray!60] (axis cs:-0.088,-0.044) -- (axis cs:-0.041,-0.071); \draw[cugray!60] (axis cs:0.181,-0.151) -- (axis cs:-0.020,-0.035); \draw[cugray!60] (axis cs:1.119,-0.188) -- (axis cs:0.198,0.344); \draw[cugray!60] (axis cs:-1.275,-1.490) -- (axis cs:-0.964,-1.670); \draw[cugray!60] (axis cs:1.485,1.146) -- (axis cs:0.867,1.502);
\addplot[only marks, mark=*, mark size=1.1pt, cuamber] coordinates {(-1.099,-1.904)(-1.098,-1.901)(-0.635,-1.100)(-0.138,-0.239)(-1.136,-1.967)(-0.034,-0.059)(-0.502,-0.869)(0.419,0.727)(0.505,0.875)(0.794,1.375)(0.804,1.392)(-0.174,-0.301)(0.466,0.807)(1.146,1.986)(-0.422,-0.730)(0.472,0.818)(0.320,0.555)(1.060,1.836)(-0.672,-1.165)(-0.524,-0.908)(0.544,0.943)(0.249,0.432)(-0.861,-1.491)(0.059,0.103)(0.414,0.718)(-0.041,-0.071)(-0.020,-0.035)(0.198,0.344)(-0.964,-1.670)(0.867,1.502)};}%
\only<6>{\addplot[cuamber, very thick, domain=-3:3, samples=2] ({x*0.0000},{x*1.0000});
\draw[cugray!60] (axis cs:-1.726,-1.542) -- (axis cs:-0.000,-1.542); \draw[cugray!60] (axis cs:-1.799,-1.496) -- (axis cs:-0.000,-1.496); \draw[cugray!60] (axis cs:-1.138,-0.810) -- (axis cs:-0.000,-0.810); \draw[cugray!60] (axis cs:0.261,-0.470) -- (axis cs:-0.000,-0.470); \draw[cugray!60] (axis cs:-2.136,-1.390) -- (axis cs:-0.000,-1.390); \draw[cugray!60] (axis cs:0.438,-0.332) -- (axis cs:-0.000,-0.332); \draw[cugray!60] (axis cs:-0.255,-1.012) -- (axis cs:-0.000,-1.012); \draw[cugray!60] (axis cs:0.653,0.592) -- (axis cs:0.000,0.592); \draw[cugray!60] (axis cs:0.203,1.050) -- (axis cs:0.000,1.050); \draw[cugray!60] (axis cs:0.986,1.264) -- (axis cs:0.000,1.264); \draw[cugray!60] (axis cs:1.187,1.170) -- (axis cs:0.000,1.170); \draw[cugray!60] (axis cs:-0.188,-0.293) -- (axis cs:-0.000,-0.293); \draw[cugray!60] (axis cs:0.417,0.835) -- (axis cs:0.000,0.835); \draw[cugray!60] (axis cs:1.056,2.038) -- (axis cs:0.000,2.038); \draw[cugray!60] (axis cs:-0.473,-0.701) -- (axis cs:-0.000,-0.701); \draw[cugray!60] (axis cs:0.661,0.709) -- (axis cs:0.000,0.709); \draw[cugray!60] (axis cs:0.355,0.534) -- (axis cs:0.000,0.534); \draw[cugray!60] (axis cs:1.733,1.447) -- (axis cs:0.000,1.447); \draw[cugray!60] (axis cs:-0.933,-1.014) -- (axis cs:-0.000,-1.014); \draw[cugray!60] (axis cs:-0.204,-1.093) -- (axis cs:-0.000,-1.093); \draw[cugray!60] (axis cs:0.087,1.207) -- (axis cs:0.000,1.207); \draw[cugray!60] (axis cs:0.374,0.360) -- (axis cs:0.000,0.360); \draw[cugray!60] (axis cs:-1.728,-0.991) -- (axis cs:-0.000,-0.991); \draw[cugray!60] (axis cs:0.064,0.100) -- (axis cs:0.000,0.100); \draw[cugray!60] (axis cs:0.683,0.563) -- (axis cs:0.000,0.563); \draw[cugray!60] (axis cs:-0.088,-0.044) -- (axis cs:-0.000,-0.044); \draw[cugray!60] (axis cs:0.181,-0.151) -- (axis cs:-0.000,-0.151); \draw[cugray!60] (axis cs:1.119,-0.188) -- (axis cs:-0.000,-0.188); \draw[cugray!60] (axis cs:-1.275,-1.490) -- (axis cs:-0.000,-1.490); \draw[cugray!60] (axis cs:1.485,1.146) -- (axis cs:0.000,1.146);
\addplot[only marks, mark=*, mark size=1.1pt, cuamber] coordinates {(-0.000,-1.542)(-0.000,-1.496)(-0.000,-0.810)(-0.000,-0.470)(-0.000,-1.390)(-0.000,-0.332)(-0.000,-1.012)(0.000,0.592)(0.000,1.050)(0.000,1.264)(0.000,1.170)(-0.000,-0.293)(0.000,0.835)(0.000,2.038)(-0.000,-0.701)(0.000,0.709)(0.000,0.534)(0.000,1.447)(-0.000,-1.014)(-0.000,-1.093)(0.000,1.207)(0.000,0.360)(-0.000,-0.991)(0.000,0.100)(0.000,0.563)(-0.000,-0.044)(-0.000,-0.151)(-0.000,-0.188)(-0.000,-1.490)(0.000,1.146)};}%
\only<7>{\addplot[cuamber, very thick, domain=-3:3, samples=2] ({x*-0.7071},{x*0.7071});
\draw[cugray!60] (axis cs:-1.726,-1.542) -- (axis cs:-0.092,0.092); \draw[cugray!60] (axis cs:-1.799,-1.496) -- (axis cs:-0.151,0.151); \draw[cugray!60] (axis cs:-1.138,-0.810) -- (axis cs:-0.164,0.164); \draw[cugray!60] (axis cs:0.261,-0.470) -- (axis cs:0.366,-0.366); \draw[cugray!60] (axis cs:-2.136,-1.390) -- (axis cs:-0.373,0.373); \draw[cugray!60] (axis cs:0.438,-0.332) -- (axis cs:0.385,-0.385); \draw[cugray!60] (axis cs:-0.255,-1.012) -- (axis cs:0.379,-0.379); \draw[cugray!60] (axis cs:0.653,0.592) -- (axis cs:0.031,-0.031); \draw[cugray!60] (axis cs:0.203,1.050) -- (axis cs:-0.423,0.423); \draw[cugray!60] (axis cs:0.986,1.264) -- (axis cs:-0.139,0.139); \draw[cugray!60] (axis cs:1.187,1.170) -- (axis cs:0.009,-0.009); \draw[cugray!60] (axis cs:-0.188,-0.293) -- (axis cs:0.052,-0.052); \draw[cugray!60] (axis cs:0.417,0.835) -- (axis cs:-0.209,0.209); \draw[cugray!60] (axis cs:1.056,2.038) -- (axis cs:-0.491,0.491); \draw[cugray!60] (axis cs:-0.473,-0.701) -- (axis cs:0.114,-0.114); \draw[cugray!60] (axis cs:0.661,0.709) -- (axis cs:-0.024,0.024); \draw[cugray!60] (axis cs:0.355,0.534) -- (axis cs:-0.090,0.090); \draw[cugray!60] (axis cs:1.733,1.447) -- (axis cs:0.143,-0.143); \draw[cugray!60] (axis cs:-0.933,-1.014) -- (axis cs:0.040,-0.040); \draw[cugray!60] (axis cs:-0.204,-1.093) -- (axis cs:0.444,-0.444); \draw[cugray!60] (axis cs:0.087,1.207) -- (axis cs:-0.560,0.560); \draw[cugray!60] (axis cs:0.374,0.360) -- (axis cs:0.007,-0.007); \draw[cugray!60] (axis cs:-1.728,-0.991) -- (axis cs:-0.369,0.369); \draw[cugray!60] (axis cs:0.064,0.100) -- (axis cs:-0.018,0.018); \draw[cugray!60] (axis cs:0.683,0.563) -- (axis cs:0.060,-0.060); \draw[cugray!60] (axis cs:-0.088,-0.044) -- (axis cs:-0.022,0.022); \draw[cugray!60] (axis cs:0.181,-0.151) -- (axis cs:0.166,-0.166); \draw[cugray!60] (axis cs:1.119,-0.188) -- (axis cs:0.653,-0.653); \draw[cugray!60] (axis cs:-1.275,-1.490) -- (axis cs:0.107,-0.107); \draw[cugray!60] (axis cs:1.485,1.146) -- (axis cs:0.170,-0.170);
\addplot[only marks, mark=*, mark size=1.1pt, cuamber] coordinates {(-0.092,0.092)(-0.151,0.151)(-0.164,0.164)(0.366,-0.366)(-0.373,0.373)(0.385,-0.385)(0.379,-0.379)(0.031,-0.031)(-0.423,0.423)(-0.139,0.139)(0.009,-0.009)(0.052,-0.052)(-0.209,0.209)(-0.491,0.491)(0.114,-0.114)(-0.024,0.024)(-0.090,0.090)(0.143,-0.143)(0.040,-0.040)(0.444,-0.444)(-0.560,0.560)(0.007,-0.007)(-0.369,0.369)(-0.018,0.018)(0.060,-0.060)(-0.022,0.022)(0.166,-0.166)(0.653,-0.653)(0.107,-0.107)(0.170,-0.170)};}%
        \end{axis}
      \end{tikzpicture}
    \end{column}
    \begin{column}{0.44\textwidth}
      \footnotesize
      \begin{block}{Project onto a line}%
        \noindent\begin{minipage}[t][6ex][t]{\dimexpr\linewidth-7pt\relax}
\only<1>{candidate axis: angle $0^\circ$, variance of projections $= 1.000$}%
\only<2>{candidate axis: angle $15^\circ$, variance of projections $= 1.423$}%
\only<3>{candidate axis: angle $30^\circ$, variance of projections $= 1.733$}%
\only<4>{\textbf{first principal component}: angle $45^\circ$, variance of projections $= 1.847$}%
\only<5>{candidate axis: angle $60^\circ$, variance of projections $= 1.733$}%
\only<6>{candidate axis: angle $90^\circ$, variance of projections $= 1.000$}%
\only<7>{candidate axis: angle $135^\circ$, variance of projections $= 0.153$}%
        \end{minipage}%
      \end{block}
      \centering
      \begin{tikzpicture}
        \begin{axis}[width=5.8cm,height=3.5cm,xmin=0,xmax=180,ymin=0,ymax=2.05,
          xlabel={axis angle (degrees)},ylabel={variance},xtick={0,45,90,135,180},
          tick label style={font=\tiny},label style={font=\scriptsize}]
          \addplot[cugray, thick] coordinates {(0,1.0000)(5,1.1470)(10,1.2896)(15,1.4233)(20,1.5442)(25,1.6486)(30,1.7332)(35,1.7956)(40,1.8338)(45,1.8467)(50,1.8338)(55,1.7956)(60,1.7332)(65,1.6486)(70,1.5442)(75,1.4233)(80,1.2896)(85,1.1470)(90,1.0000)(95,0.8530)(100,0.7104)(105,0.5767)(110,0.4558)(115,0.3514)(120,0.2668)(125,0.2044)(130,0.1662)(135,0.1533)(140,0.1662)(145,0.2044)(150,0.2668)(155,0.3514)(160,0.4558)(165,0.5767)(170,0.7104)(175,0.8530)(180,1.0000)};
\only<1>{\addplot[only marks, mark=*, mark size=2.2pt, cuamber] coordinates {(0.0,1.000)};}%
\only<2>{\addplot[only marks, mark=*, mark size=2.2pt, cuamber] coordinates {(15.0,1.423)};}%
\only<3>{\addplot[only marks, mark=*, mark size=2.2pt, cuamber] coordinates {(30.0,1.733)};}%
\only<4>{\addplot[only marks, mark=*, mark size=2.2pt, cuteal] coordinates {(45.0,1.847)};}%
\only<5>{\addplot[only marks, mark=*, mark size=2.2pt, cuamber] coordinates {(60.0,1.733)};}%
\only<6>{\addplot[only marks, mark=*, mark size=2.2pt, cuamber] coordinates {(90.0,1.000)};}%
\only<7>{\addplot[only marks, mark=*, mark size=2.2pt, cuamber] coordinates {(135.0,0.153)};}%
        \end{axis}
      \end{tikzpicture}\par
      {\raggedright\scriptsize 30 points, two correlated features, standardized. Each
      grey segment joins a point to its projection on the axis.\par}
    \end{column}
  \end{columns}
\end{frame}

\begin{frame}{What the search found}
  \begin{columns}[T,onlytextwidth]
    \begin{column}{0.5\textwidth}
      \footnotesize
      \begin{block}{The first principal component}
        The variance of the projections is largest, $1.847$, at $45^\circ$ --- the
        direction along which the points are most spread. It is the first principal
        component, $\mathbf{u}_1 = (0.7071,\ 0.7071)$. The perpendicular axis, at $135^\circ$, keeps
        only $0.153$.
      \end{block}
      \begin{block}{Variance explained}
        The total variance is $2$ (two standardized features). The first component keeps
        $1.847/2 = 92.3\%$ of it, the second $7.7\%$. One number per point
        instead of two loses less than a tenth of the information.
      \end{block}
    \end{column}
    \begin{column}{0.46\textwidth}
      \footnotesize
      \begin{block}{Worked example: one point}
        $\vx = (-1.726,\ -1.542)$. Its coordinate on the new axis:
        \[ z = \vx\cdot\mathbf{u}_1 = -2.311 \]
        \textbf{Reconstruction} --- the best approximation of $\vx$ from $z$ alone:
        \[ \hat\vx = z\,\mathbf{u}_1 = (-1.634,\ -1.634) \]
      \end{block}
      \begin{keyidea}[PCA is not linear regression]
        Regression minimises \emph{vertical} distances to predict $x_2$ from $x_1$. PCA
        minimises \emph{perpendicular} distances --- the grey segments --- and treats all
        features alike. Average squared reconstruction error here: $0.153$, exactly
        the variance left in the second component.
      \end{keyidea}
    \end{column}
  \end{columns}
\end{frame}

\begin{frame}{The PCA algorithm}
  \begin{columns}[T,onlytextwidth]
    \begin{column}{0.5\textwidth}
      \begin{defnbox}[PCA, step by step]
        \small
        \begin{enumerate}
          \item \textbf{Centre}: subtract each feature's mean. (\textbf{Scale} too, if the
                features have different units.)
          \item Compute the covariance matrix
                $\Sigma = \frac1m\sum_i \vx^{(i)}\vx^{(i)\mathsf T}$.
          \item Find its eigenvectors $\mathbf{u}_1,\mathbf{u}_2,\dots$, ordered by
                eigenvalue $\lambda_1 \ge \lambda_2 \ge \cdots$.
          \item Project: $z_j = \vx\cdot\mathbf{u}_j$ for $j = 1,\dots,k$.
        \end{enumerate}
        $\lambda_j$ is the variance captured by component $j$.
      \end{defnbox}
    \end{column}
    \begin{column}{0.46\textwidth}
      \footnotesize
      \begin{block}{Our example}
        Eigenvalues of $\Sigma$: $1.847$ and $0.153$. They are exactly the two variances the
        animation found at $45^\circ$ and $135^\circ$: the search over angles and the
        eigen-decomposition give the same answer, but the second needs no search.
      \end{block}
      \begin{pitfall}[Forgetting to centre or scale]
        \footnotesize Without centring, the first ``component'' points at the mean rather
        than along the spread. Without scaling, a feature measured in millimetres
        dominates one measured in metres simply because its numbers are bigger.
      \end{pitfall}
    \end{column}
  \end{columns}
\end{frame}

\begin{frame}{PCA on handwritten digits: 64 features to 2}
  \begin{columns}[T,onlytextwidth]
    \begin{column}{0.5\textwidth}
      \centering
      \begin{tikzpicture}
        \begin{axis}[width=6.2cm,height=5.4cm,xlabel={$z_1$},ylabel={$z_2$},
          tick label style={font=\tiny},label style={font=\scriptsize},
          legend style={font=\tiny,at={(0.02,0.98)},anchor=north west},legend cell align=left]
          \addplot[only marks, mark=*, mark size=1.3pt, cublue] coordinates {(-1.557,0.171)(-1.678,-0.264)(-1.190,0.016)(-1.758,-0.231)(-1.543,0.088)(-1.680,0.017)(-1.479,-0.159)(-1.796,0.015)(-1.258,-0.350)(-1.791,0.368)(-1.560,0.280)(-1.790,-0.218)(-2.056,-0.348)(-1.528,-0.291)(-0.810,0.294)(-1.954,-0.275)(-1.846,0.073)(-1.681,0.107)(-1.551,0.578)(-1.770,0.443)(-1.697,-0.189)(-1.573,-0.224)(-1.642,-0.241)(-1.668,-0.650)(-1.815,-0.414)(-0.838,0.786)(-1.775,-0.098)(-1.648,0.421)(-1.826,-0.336)(-2.034,-0.186)(-2.050,-0.017)(-1.679,0.256)(-2.092,0.093)(-1.403,0.117)(-1.726,-0.367)(-1.781,-0.243)(-1.909,0.091)(-1.848,-0.531)(-1.786,-0.184)(-1.738,0.228)};
          \addlegendentry{digit 0}
          \addplot[only marks, mark=square*, mark size=1.3pt, curose] coordinates {(1.387,0.658)(0.491,1.758)(1.316,1.228)(1.458,1.073)(0.657,1.783)(1.560,1.024)(0.294,1.646)(1.206,0.663)(0.445,0.944)(0.920,1.313)(0.240,1.695)(1.042,0.783)(1.178,0.893)(1.247,0.862)(1.273,1.026)(1.365,0.856)(0.530,1.099)(0.556,1.245)(0.889,1.004)(1.504,0.703)(0.692,0.699)(0.530,1.425)(0.210,1.751)(0.653,0.751)(0.285,1.848)(1.501,1.198)(0.012,0.602)(1.105,0.968)(-0.506,0.438)(-0.766,0.164)(0.423,1.313)(-0.079,0.306)(0.045,0.952)(1.481,0.627)(0.012,0.306)(0.848,0.744)(1.064,0.995)(1.454,1.131)(0.018,0.827)(0.991,0.835)};
          \addlegendentry{digit 1}
          \addplot[only marks, mark=triangle*, mark size=1.6pt, cuteal] coordinates {(0.749,-0.473)(0.854,-1.002)(0.651,-1.455)(0.888,-0.458)(0.580,-1.371)(0.574,-1.000)(0.753,-0.922)(0.928,-1.310)(1.060,-1.135)(1.246,-0.687)(1.036,-1.271)(0.882,-0.689)(0.825,-1.315)(0.584,-0.940)(1.198,-0.195)(0.965,-0.292)(0.876,-1.103)(0.256,-0.919)(0.877,-1.525)(0.705,-1.314)(0.351,-0.723)(1.068,-0.686)(0.891,-0.805)(0.829,-0.742)(1.188,-1.263)(0.212,-1.009)(1.002,-1.012)(0.968,-1.011)(0.784,-1.341)(0.387,-0.585)(1.337,-0.092)(1.127,-1.058)(0.595,-0.762)(1.116,-0.502)(1.344,-0.564)(0.998,-0.535)(0.531,-1.106)(0.919,-0.863)(0.438,-1.136)(0.703,-0.539)};
          \addlegendentry{digit 7}
        \end{axis}
      \end{tikzpicture}\par
      {\scriptsize 40 of each digit shown; PCA fitted on all 539 images of 0, 1 and 7}
    \end{column}
    \begin{column}{0.46\textwidth}
      \centering
      \begin{tikzpicture}
        \begin{axis}[width=5.8cm,height=3.6cm,xmin=0,xmax=21,ymin=0,ymax=1,
          xlabel={number of components $k$},ylabel={variance kept},
          tick label style={font=\tiny},label style={font=\scriptsize}]
          \addplot[cuamber, very thick, mark=*, mark size=1pt] coordinates {(1,0.3098)(2,0.4727)(3,0.5944)(4,0.6564)(5,0.7078)(6,0.7403)(7,0.7715)(8,0.7992)(9,0.8196)(10,0.8387)(11,0.8566)(12,0.8714)(13,0.8843)(14,0.8949)(15,0.9043)(16,0.9126)(17,0.9195)(18,0.9261)(19,0.9323)(20,0.9378)};
          \addplot[cugray, dashed, domain=0:21, samples=2] {0.9};
        \end{axis}
      \end{tikzpicture}\par
      \footnotesize\raggedright
      Each $8\times8$ image is a point in $\R^{64}$. The first two components keep
      $47.3\%$ of the variance --- yet the three digits already form three clear groups,
      found without labels. Keeping $90\%$ needs $15$ components; $95\%$ needs
      $23$.
    \end{column}
  \end{columns}
\end{frame}

\begin{frame}[fragile]{PCA in code, and when not to use it}
  \begin{columns}[T,onlytextwidth]
    \begin{column}{0.5\textwidth}
\begin{lstlisting}
from sklearn.decomposition import PCA
from sklearn.preprocessing import StandardScaler

Xs  = StandardScaler().fit_transform(X)
pca = PCA(n_components=2).fit(Xs)
pca.explained_variance_ratio_   # e.g. [0.92, 0.08]
Z     = pca.transform(Xs)            # (m, 2)
X_hat = pca.inverse_transform(Z)     # (m, n)
\end{lstlisting}
    \end{column}
    \begin{column}{0.46\textwidth}
      \footnotesize
      \begin{pitfall}[PCA is not a cure for overfitting]
        \footnotesize Reducing features with PCA to avoid overfitting throws away
        variance without asking whether it predicts $y$ --- PCA never looks at the
        labels. Use regularization (Week 3) instead.
      \end{pitfall}
      \begin{block}{Good uses today}
        Visualising high-dimensional data (often as a first look before t-SNE or UMAP);
        compressing data for storage; decorrelating features. Fit the scaler and PCA on
        the training set only, then apply them to new data.
      \end{block}
    \end{column}
  \end{columns}
\end{frame}

% ===============================================================
\section{Wrap-up and Assessment}

\begin{frame}{Summary}
  \begin{columns}[T,onlytextwidth]
    \begin{column}{0.5\textwidth}
      \begin{block}{What to remember}
        \footnotesize
        \begin{itemize}
          \item Related items: nearest learned feature vectors.
          \item Collaborative filtering suffers cold start; content-based filtering
                uses side features and works from day one.
          \item Two towers: networks map user and item features to vectors whose dot
                product predicts the rating.
          \item Large catalogues: cheap high-recall retrieval, then accurate ranking.
          \item The objective a recommender optimises is an ethical choice.
          \item PCA: projections onto eigenvectors of the covariance matrix keep the most
                variance; centre and scale first.
        \end{itemize}
      \end{block}
    \end{column}
    \begin{column}{0.46\textwidth}
      \begin{block}{Formulas}
        \[ \text{related}(i) = \argmin_{k\ne i}\|\vx^{(k)} - \vx^{(i)}\|^2 \]
        \[ \hat y = \mathbf{v}_u\cdot\mathbf{v}_m, \qquad z_j = \vx\cdot\mathbf{u}_j \]
        \[ \text{variance kept} = \frac{\lambda_1+\dots+\lambda_k}{\lambda_1+\dots+\lambda_n} \]
      \end{block}
      \begin{block}{Next}
        \footnotesize Week 9: reinforcement learning --- learning from rewards, the
        Bellman equation, and deep Q-networks.
      \end{block}
    \end{column}
  \end{columns}
\end{frame}

\begin{frame}{Descriptive questions}
  \footnotesize
  \begin{enumerate}
    \item Explain how collaborative filtering finds items related to a given item, and
          use the lecture's distance table to rank the dishes related to \emph{Pepper
          Soup}. \hfill\clo{5}
    \item Compare collaborative and content-based filtering, including how each handles
          a new item and a new user. \hfill\clo{6}
    \item Describe the two-tower architecture and its cost function. Why must the two
          towers output vectors of the same length? \hfill\clo{6}
    \item Why is it useful to normalize $\mathbf{v}_u$ and $\mathbf{v}_m$ to unit length?
          \hfill\clo{4}
    \item Explain retrieval and ranking, and the trade-off in choosing how many items to
          retrieve. \hfill\clo{5}
    \item Give two ways in which optimising engagement can harm users. \hfill\clo{4}
  \end{enumerate}
\end{frame}

\begin{frame}{Descriptive questions (continued)}
  \footnotesize
  \begin{enumerate}\setcounter{enumi}{6}
    \item State the four steps of the PCA algorithm. Why must the data be centred? When
          must it be scaled? \hfill\clo{5}
    \item For the lecture's example, explain why the variance of the projections is
          largest at $45^\circ$ and what the value $1.847$ means. \hfill\clo{5}
    \item Project $\vx = (1.0,\ 0.6)$ onto $\mathbf{u}_1 = (0.7071,\ 0.7071)$ and compute its
          reconstruction and reconstruction error. \hfill\clo{6}
    \item Distinguish PCA from linear regression, with a sketch. \hfill\clo{5}
    \item Eigenvalues $4.0, 2.5, 0.9, 0.4, 0.2$: how many components keep at least $90\%$
          of the variance? \hfill\clo{4}
    \item Why is PCA a poor remedy for overfitting? \hfill\clo{4}
  \end{enumerate}
\end{frame}

\begin{frame}{Design questions}
  \footnotesize
  \begin{enumerate}
    \item \textbf{Design a news recommender.} A national news portal wants personalised
          front pages. Design the user and article features, the two-tower model, the
          retrieval stage, and safeguards against filter bubbles and sensational content.
          \hfill\clo{12}
    \item \textbf{Design for a marketplace.} An online marketplace has $5$ million products
          and $200$ milliseconds per request. Design the retrieval and ranking pipeline,
          what is precomputed, and how you will choose the number of candidates.
          \hfill\clo{12}
    \item \textbf{Design an objective.} A video platform's recommender maximises watch time.
          Propose an alternative objective that better serves users, how you would
          measure it, and how you would test it safely. \hfill\clo{10}
    \item \textbf{Design a visualisation.} You have survey data with $40$ numeric answers
          from $2000$ respondents. Design a PCA-based visualisation: preprocessing, the
          choice of $k$, and how you would explain the axes to a non-technical audience.
          \hfill\clo{10}
  \end{enumerate}
\end{frame}

\begin{frame}{Reading}
  \begin{block}{Primary text}
    G.~James, D.~Witten, T.~Hastie and R.~Tibshirani,
    \emph{An Introduction to Statistical Learning}, 2nd ed., Springer, 2021.\\[2pt]
    \hfill{\small$\rightarrow$ \textbf{\S12.2} --- principal components analysis, including
    the proportion of variance explained and the importance of scaling;
    \textbf{\S12.3} --- matrix completion for recommender systems.}
  \end{block}
  \begin{block}{Supplementary}
    C.~M.~Bishop, \emph{Pattern Recognition and Machine Learning}, Springer, 2006.\\[2pt]
    \hfill{\small$\rightarrow$ \textbf{\S12.1} --- PCA derived as maximum-variance projection
    and as minimum-error reconstruction, the two views of this lecture.}
  \end{block}
  \vspace{1mm}
  \begin{center}
    \small\color{cugray}
    Lab (CSE 816, Week 8): collaborative filtering; deep learning for content-based
    filtering; PCA and data visualization (optional).
  \end{center}
\end{frame}

\end{document}
